%------------------------------------------------------------------------------%

\usepackage{calculo_varias_variaveis-1}

\title {Forma Polar Números Complexos}

%------------------------------------------------------------------------------%
\begin{document}

\addtitle

%------------------------------------------------------------------------------%
\section{Forma Polar ou Trigonométrica de um Nũmero Complexo}

%------------------------------------------------------------------------------%
\begin{frame}[t]{Plano Complexo}
  \begin{gather*}
    z = x + yi \quad          \in \mathbb{C} \\
    (x, y)     \hspace{7.2ex} \in \R^2
  \end{gather*}

  \addtext{9}{7,1.5}{\scalebox{0.7}[0.8]{\input{fig_complex_plane}}}

  \pause
  \begin{itemize}[<+->]
    \item Parte real       \tabto{32mm} \( \Re(z) = x \)     \\[4mm]
    \item Parte imaginária \tabto{32mm} \( \Im(z) = y \)     \\[4mm]
    \item Conjugado        \tabto{32mm} \( \bar{z} = x - yi \)        \\[4mm]
    \item Módulo           \tabto{32mm} \( \abs{z} = \rho = \sqrt{x^2+y^2\,}\) \\[4mm]
    \item Argumento        \tabto{32mm} \( \arg{(z)} = \varphi = \arctan\left(\frac{\,y\,}{x}\right)\)
  \end{itemize}
\end{frame}

%------------------------------------------------------------------------------%
\begin{frame}{Forma Polar ou Trigonométrica de um Nũmero Complexo}
  Usando
  \[
    \cos(\varphi) = \frac{a}{\rho}
    \qquad\qquad
    \pause
    a = \rho\cos(\varphi)
  \]
  \[
    \pause
    \sin(\varphi) = \frac{b}{\rho}
    \qquad\qquad
    \pause
    b = \rho\sin(\varphi)
  \]

  \pause
  \vspace{\baselineskip}
  Podemos escrever o número complexo \(z=a+bi\) como
  \[
    z = \rho \left( \cos(\varphi) + i\sin(\varphi)\right)
  \]
\end{frame}

%------------------------------------------------------------------------------%
\addtocounter{exemplo}{1}
\begin{frame}{Exemplo \theexemplo}
  Escreva os números complexos na sua forma polar
  \begin{enumerate}
    \item \(z_1 = \sqrt{3} + i\)
    \item \(z_2 = -2 + 2\sqrt{3}i\)
  \end{enumerate}
\end{frame}

\begin{frame}{Exemplo \theexemplo}
  \[
    z_1 = \sqrt{3} + i
    \qquad\qquad
    \pause
    a = \Re(z) = \sqrt{3}
    \qquad\qquad
    \pause
    b = \Im(z) = 1
  \]

  \[
    \pause
    \rho
    = \sqrt{a^2 + b^2}
    \pause
    = \sqrt{\left(\sqrt{3}\right)^2 + 1^2}
    \pause
    = \sqrt{3+1}
    \pause
    = \sqrt{4}
    \pause
    = 2
  \]

  \[
    \pause
    \sin(\varphi)
    \pause
    = \frac{b}{\rho}
    \pause
    = \frac{1}{2}
    \pause
    \qquad\qquad
    \cos(\varphi)
    \pause
    = \frac{a}{\rho}
    \pause
    = \frac{\sqrt{3}}{2}
  \]

  \[
    \pause
    \varphi
    \pause
    = \ang{30}
    \pause
    = \frac{\pi}{6} \si{\radian}
  \]

  \[
    \pause
    z_1
    \pause
    = \rho\left(\cos(\varphi)+i\sin(\varphi)\right)
    \pause
    = 2\left(\cos\left(\frac{\pi}{6}\right)+i\sin\left(\frac{\pi}{6}\right)\right)
  \]
\end{frame}

\begin{frame}{Exemplo \theexemplo}
\[
  z_2 = -2 + 2\sqrt{3}i
  \qquad\qquad
  \pause
  a = \Re(z_2) = -2
  \qquad\qquad
  \pause
  b = \Im(z_2) = 2\sqrt{3}
\]

\[
  \pause
  \rho
  = \sqrt{a^2 + b^2}
  \pause
  = \sqrt{\left(-2\right)^2 + \left(2\sqrt{3}\right)^2}
  \pause
  = \sqrt{4+4\times 3}
  \pause
  = \sqrt{16}
  \pause
  = 4
\]

\[
  \pause
  \sin(\varphi)
  \pause
  = \frac{b}{\rho}
  \pause
  = \frac{2\sqrt{3}}{4}
  \pause
  = \frac{\sqrt{3}}{2}
  \pause
  \qquad\qquad
  \cos(\varphi)
  \pause
  = \frac{a}{\rho}
  \pause
  = \frac{-2}{4}
  \pause
  = -\frac{1}{2}
\]

\[
  \pause
  \varphi
  \pause
  = \ang{120}
  \pause
  = \frac{2\pi}{3} \si{\radian}
\]

\[
  \pause
  z_2
  \pause
  = \rho\left(\cos(\varphi)+i\sin(\varphi)\right)
  \pause
  = 4\left(\cos\left(\frac{2\pi}{3}\right)+i\sin\left(\frac{2\pi}{3}\right)\right)
\]
\end{frame}

%------------------------------------------------------------------------------%
\section{Multiplicação}

%------------------------------------------------------------------------------%
\begin{frame}{Multiplicação}
  \begin{align*}
    \onslide<+->{z_1 z_2}
    \onslide<+->{& = \emph{\rho_1} \big[ \cos(\emph{\varphi_1}) + i\sin(\emph{\varphi_1})\big]}
    \onslide<+->{    \emph{\rho_2} \big[ \cos(\emph{\varphi_2}) + i\sin(\emph{\varphi_2})\big] \\[3mm]}
    \onslide<+->{& = \rho_1 \rho_2}
    \onslide<+->{  \big[ \cos(\varphi_1) + i\sin(\varphi_1)\big]
      \big[ \cos(\varphi_2) + i\sin(\varphi_2)\big] \\[3mm]}
    \onslide<+->{& = \rho_1 \rho_2
    \big[ \cos(\varphi_1)\cos(\varphi_2)
         + i\;\cos(\varphi_1)\sin(\varphi_2) \\}
    \onslide<+->{& \hspace{4.4ex}
         + i\sin(\varphi_1)\cos(\varphi_2)
       + i^2\sin(\varphi_1)\sin(\varphi_2)\big] \\[3mm]}
    \onslide<+->{& =  \rho_1 \rho_2
    \big[ \cos(\varphi_1)\cos(\varphi_2) - \sin(\varphi_1)\sin(\varphi_2) \\}
    \onslide<+->{& \hspace{4.4ex}
       + i\big(\cos(\varphi_1)\sin(\varphi_2) + \sin(\varphi_1)\cos(\varphi_2)\big)
    \big]  \\[3mm]}
    \onslide<+->{& =  \emph{\rho_1 \rho_2}
    \big[ \cos(\emph{\varphi_1+\varphi_2}) + i\sin(\emph{\varphi_1+\varphi_2}) \big]}
  \end{align*}
\end{frame}

%------------------------------------------------------------------------------%
\begin{frame}{Multiplicação}
\begin{align*}
  \onslide<+->{z_1 z_2\cdots z_n}
  \onslide<+->{& = \rho_1 \big[ \cos(\varphi_1) + i\sin(\varphi_1)\big] \\
  & \hspace{3ex} \rho_2 \big[ \cos(\varphi_2) + i\sin(\varphi_2)\big] \\
  & \hspace{4ex} \vdots \\
  & \hspace{3ex} \rho_n \big[ \cos(\varphi_n) + i\sin(\varphi_n)\big] \\[3mm]}
  \onslide<+->{& = \emph{\rho_1 \rho_2 \cdots \rho_n}}
  \onslide<+->{
      \big[
        \cos(\emph{\varphi_1+\varphi_2+\cdots+\varphi_n})
     + i\sin(\emph{\varphi_1+\varphi_2+\cdots+\varphi_n})
     \big]}
\end{align*}
\end{frame}

%------------------------------------------------------------------------------%
\addtocounter{exemplo}{1}
\begin{frame}{Exemplo \theexemplo}
  Calcule o produto dos números complexos
  \[
    z_1 = 2\left(\cos\left(\frac{\pi}{6}\right) + i\sin\left(\frac{\pi}{6}\right)\right)
  \]
  \[
    z_2 = 5\left(\cos\left(\frac{\pi}{3}\right) + i\sin\left(\frac{\pi}{3}\right)\right)
  \]
\end{frame}

\begin{frame}{Exemplo \theexemplo}
  \(z_1 = 2\left(\cos\left(\frac{\pi}{6}\right) + i\sin\left(\frac{\pi}{6}\right)\right)\) e
  \(z_2 = 5\left(\cos\left(\frac{\pi}{3}\right) + i\sin\left(\frac{\pi}{3}\right)\right)\)

  \pause
  \[
    \rho_1 = 2
    \qquad
    \pause
    \varphi_1 = \frac{\pi}{6}
  \]
  \[
    \pause
    \rho_2 = 5
    \qquad
    \pause
    \varphi_2 = \frac{\pi}{3}
  \]

  \[
    \pause
    \rho
    \pause
    = \rho_1 \rho_2
    \pause
    = 2 \times 5
    \pause
    = 10
  \]
  \[
    \pause
    \varphi
    \pause
    = \varphi_1 + \varphi_2
    \pause
    = \frac{\pi}{6} + \frac{\pi}{3}
    \pause
    = \frac{\pi + 2\pi}{6}
    \pause
    = \frac{3\pi}{6}
    \pause
    = \frac{\pi}{2}
  \]

  \[
    \pause
    z = z_1 z_2
    \pause
    = \rho\left(\cos(\varphi)+i\sin(\varphi)\right)
    \pause
    = 10\left(\cos\left(\frac{\pi}{2}\right)+i\sin\left(\frac{\pi}{2}\right)\right)
    \pause
    = 10i
  \]
\end{frame}

%------------------------------------------------------------------------------%
\addtocounter{exemplo}{1}
\begin{frame}{Exemplo \theexemplo}
  Calcule o produto dos números complexos
  \[
    z_1 = 3\left(\cos\left(\frac{2\pi}{3}\right) + i\sin\left(\frac{2\pi}{3}\right)\right)
  \]
  \[
    z_2 = \sqrt{2}\left(\cos\left(\frac{\pi}{2}\right) + i\sin\left(\frac{\pi}{2}\right)\right)
  \]
\end{frame}


\begin{frame}{Exemplo \theexemplo}
  \(z_1 = 3\left(\cos\left(\frac{2\pi}{3}\right) + i\sin\left(\frac{2\pi}{3}\right)\right)\) e
  \(z_2 = \sqrt{2}\left(\cos\left(\frac{\pi}{2}\right) + i\sin\left(\frac{\pi}{2}\right)\right)\)

  \pause
  \[
    \rho_1 = 3
    \qquad
    \pause
    \varphi_1 = \frac{2\pi}{3}
  \]
  \[
    \pause
    \rho_2 = \sqrt{2}
    \qquad
    \pause
    \varphi_2 = \frac{\pi}{2}
  \]

  \[
    \pause
    \rho
    \pause
    = \rho_1 \rho_2
    \pause
    = 3\sqrt{2}
  \]
  \[
    \pause
    \varphi
    \pause
    = \varphi_1 + \varphi_2
    \pause
    = \frac{2\pi}{3} + \frac{\pi}{2}
    \pause
    = \frac{4\pi + 3\pi}{6}
    \pause
    = \frac{7\pi}{6}
  \]

  \[
    \pause
    z = z_1 z_2
    \pause
    = \rho\left(\cos(\varphi)+i\sin(\varphi)\right)
    \pause
    = 3\sqrt{2}\left(\cos\left(\frac{7\pi}{6}\right)+i\sin\left(\frac{7\pi}{6}\right)\right)
  \]
\end{frame}

%------------------------------------------------------------------------------%
\section{Divisão}

%------------------------------------------------------------------------------%
\begin{frame}{Divisão}
  A divisão de
  \[
    z_1 = \rho_1\left(\cos(\varphi_1)+i\sin(\varphi_1)\right)
  \]
  por
  \[
    z_2 = \rho_2\left(\cos(\varphi_2)+i\sin(\varphi_2)\right)
  \]
  é
  \[
    \frac{z_1}{z_2}
    = \emph{\frac{\rho_1}{\rho_2}}
    \big[\cos(\emph{\varphi_1-\varphi_2})+i\sin(\emph{\varphi_1-\varphi_2})\big]
  \]
\end{frame}


%------------------------------------------------------------------------------%
\addtocounter{exemplo}{1}
\begin{frame}{Exemplo \theexemplo}
  Divida
  \[
    z_1 = 6\left(\cos\left(\frac{2\pi}{3}\right) + i\sin\left(\frac{2\pi}{3}\right)\right)
  \]
  por
  \[
    z_2 = 3\left(\cos\left(\frac{\pi}{2}\right) + i\sin\left(\frac{\pi}{2}\right)\right)
  \]
\end{frame}


\begin{frame}{Exemplo \theexemplo}
  \(z_1 = 6\left(\cos\left(\frac{2\pi}{3}\right) + i\sin\left(\frac{2\pi}{3}\right)\right)\) e
  \(z_2 = 3\left(\cos\left(\frac{\pi}{2}\right) + i\sin\left(\frac{\pi}{2}\right)\right)\)

  \pause
  \[
    \rho_1 = 6
    \qquad
    \pause
    \varphi_1 = \frac{2\pi}{3}
    \qquad\qquad\qquad
    \pause
    \rho_2 = 3
    \qquad
    \pause
    \varphi_2 = \frac{\pi}{2}
  \]

  \[
    \pause
    \rho
    \pause
    = \frac{\rho_1}{\rho_2}
    \pause
    = \frac{6}{3}
    \pause
    = 2
  \]
  \[
    \pause
    \varphi
    \pause
    = \varphi_1 - \varphi_2
    \pause
    = \frac{2\pi}{3} - \frac{\pi}{2}
    \pause
    = \frac{4\pi - 3\pi}{6}
    \pause
    = \frac{\pi}{6}
  \]

  \[
    \pause
    z = \frac{z_1}{z_2}
    \pause
    = \rho\left(\cos(\varphi)+i\sin(\varphi)\right)
    \pause
    = 2\left(\cos\left(\frac{\pi}{6}\right)+i\sin\left(\frac{\pi}{6}\right)\right)
  \]
  \[
    \hphantom{z = \frac{z_1}{z_2}}
    \pause
    = 2\left(\frac{\sqrt{3}}{2}\ + \frac{i}{2}\right)
    \pause
    = 3 + i
  \]
\end{frame}

%------------------------------------------------------------------------------%
\section{Lista Mínima}

%------------------------------------------------------------------------------%
\begin{frame}{Lista Mínima}
%%  Cálculo Vol. 2 do Thomas 12\textsuperscript{a} ed. -- Seção 14.8
%%  \begin{enumerate}
%%    \item Estudar o texto da seção
%%    \item Resolver os exercícios: 3, 5, 7, 12, 16, 18, 24 e 30
%%  \end{enumerate}

  \vfill
  Atenção: A prova é baseada no livro, não nas apresentações
\end{frame}

%------------------------------------------------------------------------------%
\end{document}
%------------------------------------------------------------------------------%
